% !PRINTED 180
\weberpage{180}
\weq{3}{\Sigma \xi_\nu u_\nu = \Sigma \xi'_\nu u'_\nu,}
and therefore, as in \S.\,55, (3)
\weq{4}{u'_k = \overset{i}{\Sigma} \alpha_i^{(k)} u_i.}
Now
\weq{5}{u_k = \overset{i}{\Sigma} a_{k,i} x_i, \qquad u'_k = \overset{i}{\Sigma} a'_{k,i} x'_i,}
and accordingly (4) gives us the transformation of $m$ linear
functions. The coefficients of the given system are $a'_{k,h}$,
and those of the transformed system, which one obtains by substituting
the first system (5) into (4), are
\weq{6}{b_{k,h} = \overset{i}{\Sigma} \alpha_i^{(k)} a_{i,h}.}
The determinant of the former system
\[
H' = \Sigma \pm a'_{1,1} a'_{2,2} \ldots a'_{m,m}
\]
is, by \S.\,55, (4),
\[
= r \Sigma \pm b_{1,1} b_{2,2} \ldots b_{m,m},
\]
and by (6) the determinant of the $b_{k,h}$ is
\[
= r \Sigma \pm a_{1,1} a_{2,2} \ldots a_{m,m}.
\]
If, then, we set
\weq{7}{H = \Sigma \pm a_{1,1} a_{2,2} \ldots a_{m,m}}
we have the important theorem
\weq{8}{H' = r^2 H.}
The determinant $H$ is called the \gesperrt{determinant of the form}
$\varphi(x)$.

Thus under linear transformation of the function $\varphi$ it
changes only by the factor $r^2$, and it cannot vanish for the
transformed form if it does not do so for the original one.

\gesperrt{The vanishing of the determinant $H$ has the meaning
that the function $\varphi(x)$ can be regarded as
a function of fewer than $m$ variables, which
are linear functions of the $x$.}

For if, in the first place, $\varphi(x)$ passes by a linear transformation
into the function $\Phi(x')$, which is independent of one of the variables,
say of $x'_1$, then $H'$ vanishes, since the elements
in one of the rows (and columns) all vanish, and hence
by (8) also $H = 0$.

In the second place, however, $H$ is the determinant of the system of linear
equations
%CONT>>
